\DocumentMetadata{
  pdfversion=2.0,
  pdfstandard=ua-2,
  lang=en-US,
  tagging=on,
  tagging-setup={math/setup=mathml-SE,tabsorder=structure}
}
\documentclass[11pt,letterpaper]{article}
\usepackage[margin=0.68in,includefoot]{geometry}
\usepackage{newcomputermodern}
\usepackage{amsmath,amscd,mathtools}
\usepackage{tikz,adjustbox}
\providecommand{\square}{\mdlgwhtsquare}
\usepackage[hidelinks]{hyperref}
\hypersetup{
  pdftitle={Analysis First-Year Exam, August 2025, Part 2},
  pdfauthor={University of Florida Department of Mathematics},
  pdfsubject={Graduate mathematics examination},
  pdfdisplaydoctitle=true,
  bookmarksopen=true
}
\setcounter{secnumdepth}{0}
\setlength{\parindent}{0pt}
\setlength{\parskip}{0.28em}
\setlength{\emergencystretch}{3em}
\setlength{\leftmargini}{2.15em}
\setlength{\leftmarginii}{2.65em}
\setlength{\leftmarginiii}{3em}
\renewcommand{\labelenumi}{\textbf{\theenumi.}}
\pagestyle{empty}
\raggedbottom
\allowdisplaybreaks
\newenvironment{examproblems}[1]{%
  \begin{enumerate}%
  \setcounter{enumi}{#1}%
  \setlength{\topsep}{0.35em}%
  \setlength{\partopsep}{0pt}%
  \setlength{\itemsep}{0.48em}%
  \setlength{\parsep}{0.18em}%
}{\end{enumerate}}
\newenvironment{examsubparts}{%
  \begin{enumerate}%
  \setlength{\topsep}{0.18em}%
  \setlength{\partopsep}{0pt}%
  \setlength{\itemsep}{0.16em}%
  \setlength{\parsep}{0.1em}%
}{\end{enumerate}}
\newenvironment{examinstructions}{%
  \par\begingroup\itshape
}{%
  \par\endgroup\vspace{0.18em}
}
\makeatletter
\renewcommand\section{\@startsection{section}{1}{0pt}%
  {-1.25ex plus -.25ex minus -.1ex}%
  {0.55ex plus .1ex}%
  {\normalfont\large\bfseries}}
\renewcommand{\@maketitle}{%
  \newpage\null\vskip -1em%
  {\footnotesize\bfseries
    \href{https://www.ufl.edu/}{University of Florida}, \href{https://math.ufl.edu/}{Department of Mathematics}\par}%
  \vskip -0.5em%
  \tagstructbegin{tag=Title}%
    {\LARGE\bfseries\href{https://gma.math.ufl.edu/exams/analysis-first-year/}{Analysis First-Year Exam}, August 2025, Part 2\par}%
  \tagstructend%
  \vskip 0.25em%
  \tagmcbegin{artifact}%
    \hrule height 1pt%
  \tagmcend%
  \vskip 1em%
}
\makeatother
\title{Analysis First-Year Exam, August 2025, Part 2}
\author{}
\date{}
\begin{document}
\maketitle
\thispagestyle{empty}
\begin{examinstructions}
Answer FOUR questions in detail. State carefully any results used without proof.
\end{examinstructions}
\begin{examproblems}{0}
\item
This problem has two parts.
\begin{examsubparts}
\item[\textnormal{(i)}]
Let \(f:[0,1]\to\mathbb{R}\) be Riemann integrable. Must the function~\(F\) defined on \([0,1]\) by \(F(x)=\int_0^x f(t)\,dt\) be continuous? Explain.
\item[\textnormal{(ii)}]
Let \(g:[0,1]\to\mathbb{R}\) be Lebesgue integrable. Must the function~\(G\) defined on \([0,1]\) by \(G(x)=\int_{[0,x]}g\,d\lambda\) be continuous (\(\lambda\): Lebesgue measure)? Explain.
\end{examsubparts}
\item
Equip the set \(C[0,1]\) (comprising all continuous real-valued functions on the interval \([0,1]\)) with the sup metric and let 
\[
B=\{f\in C[0,1]:\sup|f|\leq 1\}
\]
 be its closed ball of radius one about the zero function. Is~\(B\) complete? Is~\(B\) compact? Explain.
\item
Assume that the continuous function \(f:[0,1]\to\mathbb{R}\) satisfies 
\[
\int_0^1 f(t)\cos(nt)\,dt=0
\]
 whenever~\(n\) is a non-negative integer. Does it follow that~\(f\) is the zero function? Proof or counterexample required.
\item
Let \((f_n)_{n=0}^{\infty}\) be a sequence of real-valued measurable functions on~\(\Omega\). Prove that if \((f_n)_{n=0}^{\infty}\) converges pointwise to \(f:\Omega\to\mathbb{R}\), then~\(f\) is measurable.
\item
The bounded function \(f:[0,1]\times[0,1]\to\mathbb{R}\) is continuous in each variable separately: thus, \(f(x,y)\) is continuous in~\(x\) for each fixed~\(y\) and continuous in~\(y\) for each fixed~\(x\). Prove the continuity of~\(F\) defined by 
\[
F(x)=\int_0^1 f(x,y)\,dy.
\]
\end{examproblems}
\end{document}
